\UseRawInputEncoding
\chapter{Matematica elementare}


\begin{esercizio}
	Risolvere la seguente disequazione: 
	\[
	|x^2-5x+6|=x^2-3
	\]
\end{esercizio}
\textit{Soluzione.} L'equazione � equivalente alla coppia di sistemi, ottenuti studiando il segno dell'argomento del valore assoluto: 
\[
\begin{cases}
	x^2-5x+6\geq 0\\
	x^2-5x+6=x^2-3
\end{cases}
\quad \vee \qquad 
\begin{cases}
	x^2-5x+6<0\\
	-x^2+5x-6=x^2-3
\end{cases} 
\]
\begin{itemize}
	\item Si risolve il primo sistema:
	\[
	\begin{cases}
		x\leq 2 \, \, \, \vee \, \, \, x\geq 3\\
		x=9/5
	\end{cases}
	\Rightarrow \quad 
	S_1=\big \{9/5\big \}
	\]
	\item Si risolve il secondo sistema: 
	\[
	\begin{cases}
		2<x<3\\
		x=1 \, \, \, \vee \, \, \, x=3/2
	\end{cases}
	\Rightarrow \quad 
	S_2=\emptyset
	\]
\end{itemize}
L'equazione ammette quindi la sola soluzione $9/5$. 
\begin{esercizio}
	Risolvere la seguente disequazione: 
	\[
	|x-2|=x-3
	\]
\end{esercizio}
\textit{Soluzione}. Anzich� studiare il comportamento del valore assoluto, si osserva che il secondo membro eredita la propriet� di non-negativit�. \' E quindi possibile elevare entrambi i membri al quadrato: 
\[
\begin{cases}
	x-3\geq 0\\
	(x-2)^2=(x-3)^2
\end{cases}
\Rightarrow \quad 
\begin{cases}
	x\geq 3\\
	2x=5
\end{cases}
\Rightarrow \quad 
\begin{cases}
	x=5/2\\
	x\geq 3
\end{cases}
\Rightarrow \quad 
S=\emptyset
\]

\begin{esercizio}
	Risolvere l'equazione: 
	\[
	x^3+(x-\sqrt{3})^3=0
	\]
\end{esercizio}
\textit{Soluzione}. L'equazione si riconduce alla forma:
\[
x^3=-(x-\sqrt{3})^3=0 \quad \Rightarrow \quad 
x^3=(\sqrt{3}-x)^3
\] 
\' E possibile estrarre le radici cubiche di entrambi i membri:
\[
x=\sqrt{3}-x
 \quad \Rightarrow \quad 
 x=\frac{\sqrt{3}}{2}
\]

\begin{esercizio}
	Risolvere la seguente disequazione:
	\[
	(x+3)^3(4-x^2)^4(-x^2+6x-5)^5\geq 0
	\]
\end{esercizio}
\textit{Soluzione}. \ \begin{itemize} 
	\item Si osserva che $(4-x^2)^4$ � non negativo e che $x=\pm 2$ sono soluzioni. \\
	Posto $x=\pm 2$ soluzioni, � quindi possibile semplificare il fattore senza alterare il segno. 
	\item Il segno di $(x+3)^3$ � lo stesso di quello di $x+3$.
	\item Il segno di $(-x^2+6x-5)^5$ � lo stesso di quello di $-x^2+6x-5$. 
\end{itemize}
La disequazione � quindi equivalente alla seguente: 
\[
x=\pm 2 \, \, \, \vee \, \, \, (x+3)(-x^2+6x-5)\geq 0
\]
Studiando il segno dei fattori si deduce l'insieme soluzione: 
\[
x\leq -3 \, \, \, \vee \, \, \, x=-2 \, \, \, \vee \, \, \, 1\leq x\leq 5
\]

\begin{esercizio}
	Risolvere la disequazione: 
	\[
	\frac{1}{(x-1)^3}\leq -\frac{1}{x^3}
	\]
\end{esercizio}
\textit{Soluzione}. Si osserva che � errato passare a reciproco entrambi i membri perch� ci� sarebbe falso nel caso in cui i due membri hanno segno opposto. 
\' E invece possibile estrarre la radice di entrambi i membri: 
\[
\frac{1}{x-1}\leq -\frac{1}{x}
\quad \Rightarrow \quad 
\frac{2x-1}{x(x-1)}\leq 0
\quad \Rightarrow \quad
x<0  \quad \Rightarrow \quad 
\frac{1}{2}\leq x<1
\]

\begin{esercizio}
	Risolvere la seguente disequazione:
	\[
\frac{x^2+4x+2}{x+1}>\sqrt{x^2}	
	\]
\end{esercizio}
\textit{Soluzione}. Si osserva che $\sqrt{x^2}=|x|$. \\
La disequazione viene quindi risolta studiando il comportamento del valore assoluto:
\[
\begin{cases}
	x\geq 0\\[6pt]
	\dfrac{x^2+4x+2}{x+1}>x
\end{cases}
\quad \vee \qquad 
\begin{cases}
	x<0\\[6pt]
	\dfrac{x^2+4x+2}{x+1}>-x
\end{cases}
\]
\begin{itemize}
	\item Si risolve il primo sistema:
	\[
	\begin{cases}
		x\geq 0\\[6pt]
		\dfrac{3x+2}{x+1}>0
	\end{cases}
	\quad \Rightarrow \quad 
	S_1: \, x\geq 0
	\]
	Infatti, sotto la condizione $x\geq 0$, sia il numeratore sia il denominatore sono positivi. 
	\item Si risolve il secondo sistema: 
	\[
	\begin{cases}
		x<0\\[6pt]
		\dfrac{2x^2+5x+4}{x+1}>0
	\end{cases}
	\Rightarrow \quad 
	S_1: \, -2<x<-1 \, \, \, \vee \, \, \, -\frac{1}{2}<x\leq 0
	\] 
	Unendo le soluzioni: 
	\[
	S=S_1\cup S_2: \, \, -2<x<-1 \, \, \, \vee \, \, \, x>-\frac{1}{2}
	\]
\end{itemize}

\begin{esercizio}\label{equazrazionale}
	Mostrare che se un numero razionale $\overline{x}=p/q$, con $p\in \mathbb{Z}$, $q\in \mathbb{N}^*$ e $p,q$ primo tra loro) � soluzione dell'equazione algebrica: 
	\[
	p(x)=\sum_{i=0}^n a_i x_i=0
	\]
	allora $p$ � divisore del termine noto $a_0$ mentre $q$ � divisore del coefficiente direttore $a_n$. 
\end{esercizio}
\textit{Soluzione}. Per ipotesi, $\overline{x}$ � soluzione dell'equazione $p(x)=0$:
\[
	p(\overline{x})=0 
	\quad \Rightarrow \quad 
	a_n \, \frac{p^n}{q^n}+a_{n-1} \, \frac{p^{n-1}}{q^{n-1}}+\dots+a_1 \, \frac{p}{q}+a_0=0
\]
Essendo $q\neq 0$, si moltiplica per $q^n$:
\[
a_n p^n+a_{n-1} \, p^{n-1}q+\dots+a_1pq^{n-1}+a_0q^n=0 \]
\begin{itemize}
	\item Si evidenzia un fattore $p$:
	\[
	p \, \big (a_n p^{n-1}+\dots +a_1q^{n-1}\big )=-a_0q^n 
	\]
	Siccome $p$ divide $a_0q^n$ e $p$ � primo con $q$, si ha che $p$ divide $a_0$.
	\item Si evidenzia un fattore $q$:
	\[
	q \, (a_0q^{n-1}+a_1pq^{n-2}+\dots+a_{n-1}p^{n-1}=-a_np^n
	\]
	Siccome $q$ divide $a_n p^n$ e $q$ � primo con $p$, si ha che $q$ divide $a_n$. 
\end{itemize}
Si osserva che la condizione di divisibilit� � necessaria affinch� $\overline{n}\in \mathbb{Q}$ sia soluzione. Quindi se tale condizione non � soddisfatta allora $\overline{x}\in \mathbb{Q}$ non � soluzione dell'equazione.\\
Segue inoltre che gli zeri razionali di un polinomio sono del tipo $\pm a/b$, dove $a$ � un divisore del termine noto e $b$ � un divisore del coefficiente direttore. \\
Per esempio, gli zeri razionali del polinomio $p(x)=2x^3-13x^2+18x-6$ si ricercano tra i seguenti: 
\[
\pm 1 \qquad \pm 2 \qquad \pm 3 \qquad \pm 6 \qquad \pm \frac{1}{2} \qquad \pm \frac{3}{2}
\]

\begin{esercizio}
	Risolvere la seguente disequazione: 
	\[
	2^x>2+2^{-x}
	\]
\end{esercizio}
\textit{Soluzione}. Si porta la disequazione in forma normale: 
\[
2^{2x}-2\cdot 2^x-1>0
\]
Posto $2^x=t$, l'equazione associata � riconducibile a una forma di secondo grado: 
\[
t^2-2t-1=0 
\quad \Rightarrow \quad 
t=1\pm \sqrt{2}
\]
Da cui fattorizzando:
\begin{align*}
&\big (2^x-(1-\sqrt{2})\big ) \, \big (2^x-(1+\sqrt{2})\big )>0
\quad \Rightarrow \quad 
\big (2^x+\sqrt{2}-1)\big ) \, \big (2^x-(1+\sqrt{2})\big )>0\\
\end{align*}
Il primo fattore � positivo: 
\[
2^x-(1+\sqrt{2})>0 \quad \Rightarrow \quad 
2^x>1+\sqrt{2}
\quad \Rightarrow \quad 
x>\log_2 (1+\sqrt{2})
\]

\begin{esercizio}
	Risolvere la seguente equazione: 
	\[
	\sin x-\cos x=1
	\]
\end{esercizio}
\textit{Soluzione}. L'equazione si pu� risolvere applicando le formule di addizione oppure il metodo dell'angolo aggiunto. 
Facendo uso delle formule parametriche, posto $t=\tan (x/2)$:
\[
\frac{x}{2}\neq \frac{\pi}{2}+k\pi 
\quad \Rightarrow \quad 
x\neq \pi +k2\pi 
\]
Si osserva che $x=\pi +k2\pi$ � soluzione. Posto $x\neq \pi +k2\pi$:
\[
\sin x-\cos x=1
\quad \Rightarrow \quad 
\frac{2t}{1+t^2}+\frac{1-t^2}{1+t^2}=1
\quad \Rightarrow \quad 
\frac{2t+1-t^2}{1+t^2}=1
\quad \Rightarrow \quad 
t=1
\]
Si ha: 
\[
t=\tan \frac{x}{2}=1
\quad \Rightarrow \quad 
\frac{x}{2}=\frac{\pi}{4}+k\pi 
\quad \Rightarrow \quad 
x=\frac{\pi}{2}+k2\pi 
\]
L'equazione ammette quindi le seguenti soluzioni: 
\[
x=\frac{\pi}{2}+k2\pi \, \, \, \vee \, \, \, x=\pi +k2\pi 
\]

\begin{esercizio}
	Risolvere, nell'intervallo $(-2\pi, 2\pi )$ la seguente disequazione: 
	\[
	\sin 2x>\cos x
	\]
\end{esercizio}
\textit{Soluzione}. 
\[
2\sin x \cos x-\cos x>0 \quad \Rightarrow \quad 
\cos x \, (2\sin x-1)>0
\]
Anzich� studiare il segno dei fattori, si pu� procedere per casi: 
\[
\begin{cases}
	\cos x>0\\
	2\sin x-1>0
	\end{cases} \quad \vee \qquad 
	\begin{cases}
		\cos x<0\\
		2\sin x-1<0
	\end{cases}
\]
Da cui le soluzioni: 
\[
-\frac{11}{6} \, \pi <x<-\frac{3}{2} \, \pi \quad \vee \quad 
-\frac{7}{6} \, \pi <x<-\frac{\pi}{2}
 \quad \vee \quad 
\frac{\pi}{6}<x<\frac{\pi}{2}
 \quad \vee \quad 
 \frac{5}{6} \, \pi <x<\frac{3}{2} \, \pi 
\]

\begin{esercizio}
	Risolvere l'equazione: 
	\[
	\sqrt{2\log_2 x+3}=\log_2 x
	\]
\end{esercizio}
\textit{Soluzione}. 
\[
\begin{cases}
x>0\\
\log_2 x>0\\
2\log_2 x+3\geq 0\\
2\log_2 x+3=\log_2^2 x
\end{cases}
\Rightarrow \quad 
\begin{cases}
	x>0\\
	\log_2 x>0\\
	\log_2^2 x-2\log_2 x-3=0
\end{cases}
\Rightarrow \quad 
\begin{cases}
	x>0\\
	\log_2 x>0\\
	\log_2 x=-1 \, \, \, \vee \, \, \, \log_2 x=3
\end{cases}
\]
Da cui: 
\[
\log_2 x=3 \quad \Rightarrow \quad 
x=8
\]


\begin{esercizio}
	Risolvere la disequazione: 
	\[
	\sqrt{2\log_2 x+3}>\log_2 x
	\]
\end{esercizio}
\textit{Soluzione}. 
\[
\begin{cases}
	x>0\\
	\log_2 x<0\\
	2\log_2 x+3\geq 0
\end{cases}
\qquad \begin{cases}
	x>0\\
	\log_2 x\geq 0\\
	2\log_2 x+3\geq \log_2^2 x\\
\end{cases}
\]
I due sistemi ammettono rispettivamente i seguenti insiemi soluzione: 
\[
\frac{1}{2\sqrt{2}}\leq x<1 \qquad \qquad 1\leq x<2^3
\]
Da cui la soluzione della disequazione: 
\[
\frac{1}{2\sqrt{2}}\leq x<8
\]

\begin{esercizio}
	Risolvere la seguente equazione: 
	\[
	2\sqrt[3]{(e^x-1)^2}-e^x=0
	\]
\end{esercizio}
\textit{Soluzione}. \' E possibile elevare al cubo entrambi i membri: 
\[
8 \, (e^x-1)^2=e^{3x} \quad \Rightarrow \quad 
e^{3x}-8e^{2x}+16e^x+8=0
\]
Fattorizzando:
\[
(e^x-2)(e^{2x}-6e^x+4)=0
\quad \Rightarrow \quad 
x=\ln 2 \, \, \, \vee \, \, \, x=\ln (3\pm \sqrt{5})
\]

\begin{esercizio}
	Determinare le soluzioni positive della seguente equazione: 
	\[
	(1+x) \, \sqrt{x^2+2x}=2\sqrt{3}
	\]
\end{esercizio}
\textit{Soluzione}. Siccome, per ipotesi, $x>0$ la radice � definita ed entrambi i membri sono positivi. \' E quindi possibile elevare entrambi i membri al quadrato senza introdurre radici positive estranee:
\[
(1+x)^2 \, (x^2+2x)=12
\quad \Rightarrow \quad 
x^4+4x^3+5x^2+2x-12=0
\quad \Rightarrow \quad 
(x-1)(x^3+5x^2+10x+12)=0
\]
Essendo $x>0$, il secondo fattore � positivo. Si conclude che $x=1$ � l'unica soluzione positiva. 

\begin{esercizio}
	Risolvere la disequazione: 
	\[
	\sqrt{\sin x}\leq 1-\cos x
	\]
\end{esercizio}
\textit{Soluzione}. Il secondo membro � non negativo. Dopo aver posto la condizione di esistenza del radicale, � possibile elevare entrambi i membri al quadrato: 
\[
\begin{cases} 
\sin x\geq 0\\
(1-\cos x)^2\geq \sin x
\end{cases} 
\]
Tenendo presente che $\sin x\geq 0$, si pu� esprimere in maniera univoca $\sin x$ in funzione di $\cos x$:
\[
\sin x\geq 0 \quad \Rightarrow \quad 
\sin x=\sqrt{1-\cos^2 x}
\quad \Rightarrow \quad 
\begin{cases}
\sin x\geq 0\\
(1-\cos x)^2\geq \sqrt{1-\cos^2 x}
\end{cases}
\]
\' E possibile elevare al quadrato: 
\[
(1-\cos x)^4\geq 1-\cos2 x \quad \Rightarrow \quad 
(1-\cos x)^4=(1-\cos x)(1+\cos x)
\]
Dopo aver osservato che $x=k2\pi $ � soluzione: 
\[
(1-\cos x)^3=1+\cos x
\quad \Rightarrow \quad 
\cos^3x-3\cos^2 x+4\cos x\leq 0
\quad \Rightarrow \quad 
\cos x \, (\cos^2 x-3\cos x+4)\leq 0
\]
Siccome il trinomio � ovunque positivo, la disequazione si riduce alla forma seguente: 
\[
\begin{cases} 
	\sin x\geq 0\\
	\cos x\leq 0
	\end{cases} 
\Rightarrow \quad 
S: \, \frac{\pi}{2}+k2\pi <x<\pi+k2\pi \quad \vee \quad x=k2\pi 
\]